\section{Literature Review}
\label{sec:literature_review}

The literature review first considers research on phonological learning and
task-trained neural organization (Sections~\ref{sec:related} and
\ref{sec:lr-task-trained}). It then examines Mandarin data sources and neural
representations (Sections~\ref{sec:lr-data} and \ref{sec:lr-models}),
followed by phonetic supervision and representation evaluation
(Sections~\ref{sec:lr-supervision} and \ref{sec:lr-evaluation}).

\subsection{Related Works}
\label{sec:related}

First works on computational phonetics have investigated how neural networks
acquire pronunciation knowledge from written language. Seidenberg and McClelland
developed a distributed model that learns mappings between orthographic and
phonological representations through supervised training
\cite{ref_seidenberg1989}. Plaut et al. extended this approach to explain
regular and irregular word reading using learned orthography-to-phonology
mappings \cite{ref_plaut1996}. These studies demonstrated that neural networks
can acquire systematic pronunciation relationships without relying on explicit
pronunciation rules. Their evaluations focused primarily on reading performance
and the ability to reproduce observed human reading effects.

Computational phonetic learning has also been studied using representations
derived directly from speech. Bartelds et al. examined speech representations
from pretrained neural networks and evaluated their ability to capture
pronunciation variation \cite{ref_bartelds2022}. Distances between the learned
representations showed correspondence with human perceptual judgments.
Matusevych et al. compared computational models of infant phonetic learning
based on different learning mechanisms \cite{ref_matusevych2023}. The models
were evaluated through phonetic discrimination tasks using speech data from
different languages. These studies demonstrate that continuous speech
representations can capture phonetic relationships and that different learning
mechanisms may affect their organization. However, the models operate on
acoustic inputs rather than written character sequences.

Chinese character reading provides a setting for investigating the acquisition
of phonological information from structured orthographic inputs. Xing et al.
developed a self-organizing connectionist model of Chinese character acquisition
\cite{ref_xing2002}. The model reproduced effects of character frequency and
phonological regularity observed in reading. Yang et al. subsequently developed
a triangle model of Chinese reading that represented relationships between
orthography, phonology, and semantics \cite{ref_yang2006}. Their later
statistical learning model predicted distributed phonological feature vectors
from Chinese orthography \cite{ref_yang2009}. The model reproduced
character-naming effects associated with frequency, regularity, and consistency.
Its phonological targets encoded relationships between speech sounds through
shared features. However, the evaluation primarily concerned pronunciation
performance and reading effects rather than independent graded similarity in
hidden representations. Reichle and Yu reviewed these computational approaches
and examined their assumptions about Chinese reading \cite{ref_reichle2018}.

More recent studies have examined phonological representations and training
conditions in Chinese reading. Li et al. used representational similarity
analysis to investigate how human brain activity reflects phonological
information during Chinese character reading \cite{ref_li_chinese2022}. Their
evaluation compared neural response patterns with phonological dissimilarity
measures. Kuo and Chang investigated the effects of reading instruction using
computational triangle models of Chinese character reading \cite{ref_kuo2025}.
Their models were trained under sound-focused, meaning-focused, and balanced
instructional conditions. The authors evaluated pronunciation performance,
reading effects, and the relative reliance on semantic and phonological
pathways. The training conditions varied the emphasis on reading tasks rather
than the form of phonetic information provided as supervision. Furthermore, the
evaluation did not directly examine correspondence between hidden-state
similarities and independent graded phonetic distances.

Previous studies have demonstrated that neural networks can learn phonological
information from orthography and that different learning objectives can
influence neural organization. Other studies have examined phonetic
relationships in speech representations and human neural responses. However, how
different forms of phonetic supervision influence graded similarity in
contextualized character representations remains insufficiently understood. We
investigate this question by comparing four training objectives: complete Pinyin
syllable prediction, phonological component prediction, continuous acoustic
alignment, and combined component-based and acoustic supervision. We evaluate
the resulting representation spaces through their correspondence with
independent, perceptually informed Mandarin phonetic distances.


\subsection{Task-Trained Networks and Internal Representations}
\label{sec:lr-task-trained}

The relationship between learning objectives and internal neural organization
has also been examined in computational neuroscience. Sussillo et al. trained
recurrent neural networks to generate muscle activity and identified dynamics
resembling aspects of biological motor processing \cite{ref_sussillo2015}. Kell
et al. demonstrated that task-optimized networks for speech and music
recognition could predict aspects of human auditory cortical responses
\cite{ref_kell2018}. Correspondences between speech recognition models and
neural activity during auditory processing were further examined by Li et al.
and Brodbeck et al. \cite{ref_li_auditory2023,ref_brodbeck2025}. Wu et al.
showed that multitask learning could promote modular organization in recurrent
networks \cite{ref_wu2026}. Together, these findings support the investigation
of internal representations as outcomes of neural network training, without
assuming that artificial networks reproduce the mechanisms of biological
learning.


\subsection{Data Sources for Mandarin Pronunciation Learning}
\label{sec:lr-data}

\subsubsection{Character--Pinyin Corpora}
\label{sec:lr-corpora}

Mandarin pronunciation prediction requires character sequences with
corresponding pronunciation annotations. The Chinese Polyphone with Pinyin (CPP)
benchmark introduced for g2pM provides annotated sentences for evaluating
context-dependent pronunciation prediction \cite{ref_g2pm2020}. The same
benchmark has been used to assess subsequent neural approaches to polyphone
disambiguation \cite{ref_g2pw2022}. Beyond polyphonic characters, learning
representations across a wider pronunciation inventory requires annotations for
character occurrences throughout the corpus. The Duyu/Pinyin-Hanzi dataset
provides Chinese text sequences paired with Pinyin transcriptions and represents
a possible source for this purpose \cite{ref_duyu}. Additionally, the AISHELL-3
corpus provides Mandarin text with Pinyin annotations \cite{ref_aishell32021}.
These resources support the training of neural models to predict pronunciations
from written characters, either as complete syllables or as phonological
components.

\subsubsection{Mandarin Speech Recordings}
\label{sec:lr-speech-recordings}

Acoustic representations require recorded pronunciations from which speech
features can be extracted. AISHELL-3 provides Mandarin speech recordings
accompanied by character-level and Pinyin-level transcriptions
\cite{ref_aishell32021}. Such corpora support the analysis of pronunciation in
continuous speech but require segmentation or alignment to obtain individual
syllable recordings. Audio-assisted pronunciation modeling has addressed this
problem through alignment between speech and written transcripts
\cite{ref_g2pa2024}. Alternatively, collections of isolated syllable
pronunciations provide direct access to recordings associated with individual
Pinyin labels. The Chinese Lessons sound collection is an example of such a
resource \cite{ref_chinese_lessons}. The choice between continuous speech and
isolated syllables determines whether the extracted features describe contextual
acoustic realizations or reference pronunciations.


\subsection{Neural Models and Phonetic Representations}
\label{sec:lr-models}

\subsubsection{Character Encoding and Contextual Models}
\label{sec:lr-character-models}

Neural models of Mandarin pronunciation must represent individual characters
while accounting for their linguistic context. Earlier Chinese reading models
incorporated structured orthographic descriptions, including character
components, to predict pronunciation. More recent grapheme-to-phoneme models
have used learned character embeddings and contextual sequence encoders. Park
and Lee employed a bidirectional recurrent network for Mandarin polyphone
disambiguation \cite{ref_g2pm2020}. Pretrained Transformer encoders provide
another approach to contextual character representation, as demonstrated by BERT
and its application in g2pW \cite{ref_bert2019,ref_g2pw2022}. These
architectures offer different mechanisms for contextualization, while their
hidden states can be examined separately from the pronunciation predictions
produced by their output layers.

\subsubsection{Representations of Phonetic Information}
\label{sec:lr-phonetic-representations}

Phonetic information can be expressed through symbolic categories or continuous
feature representations. Distributed phonological vectors encode properties
shared by different pronunciations and have been used as targets in
computational models of Chinese reading \cite{ref_yang2009,ref_kuo2025}. Pinyin
provides a symbolic description of Mandarin syllables, including their segmental
composition and lexical tones. It has been used both for pronunciation
prediction and as additional input information in Chinese language models
\cite{ref_g2pm2020,ref_chinesebert2021}. The latter approach differs from
learning pronunciation as an output target, since phonetic information is
explicitly supplied to the encoder. Continuous acoustic representations provide
another description of pronunciation, with features learned from recorded speech
rather than predefined phonological categories \cite{ref_bartelds2022}. These
alternatives differ in the information they encode and in the relationships they
can express between distinct pronunciations.


\subsection{Phonetic Supervision and Learning Objectives}
\label{sec:lr-supervision}

\subsubsection{Categorical Pronunciation Prediction}
\label{sec:lr-categorical}

Pronunciation prediction commonly treats possible pronunciations as discrete
classes. Models such as g2pM and g2pW learn to select a contextually appropriate
pronunciation from a set of candidates \cite{ref_g2pm2020,ref_g2pw2022}. An
alternative is to decompose a syllable into phonological components and predict
them separately. Li et al. examined the granularity of pronunciation prediction
as an auxiliary task for Chinese spelling correction \cite{ref_li_scope2022}.
Liang et al. similarly investigated disentangled phonetic information for
correcting Chinese characters \cite{ref_liang2023}. These approaches demonstrate
that pronunciation can be represented through either complete categories or
smaller phonological units. Their primary evaluations, however, concern task
performance rather than the graded geometry of contextual hidden
representations.

\subsubsection{Acoustic Representation Learning and Alignment}
\label{sec:lr-acoustic-alignment}

Self-supervised speech encoders provide a means of extracting continuous
features from raw audio. The wav2vec 2.0 framework learns speech representations
through a contrastive objective, while XLSR extends this approach to
multilingual speech \cite{ref_wav2vec2,ref_xlsr2021}. Liao et al. employed a
Mandarin-fine-tuned XLSR-53 model to extract acoustic representations from
Chinese speech recordings \cite{ref_liao2026}. The extracted features were used
in a multimodal system for detecting mild cognitive impairment, demonstrating
the application of this model as an acoustic feature extractor.

Speech-derived information has also been incorporated into models operating
across acoustic and textual modalities. Chen et al. proposed MAESTRO to
learn matched representations from speech and text \cite{ref_maestro2022}. Yang
used aligned speech recordings to improve Mandarin pronunciation disambiguation
\cite{ref_g2pa2024}. Mazumdar et al. introduced a contrastive approach that
aligns written graphemes with speech-derived representations
\cite{ref_sparcle2026}. These studies establish approaches for combining
symbolic linguistic information with acoustic features. Nevertheless, their
objectives primarily concern speech recognition, pronunciation prediction, or
synthesis, rather than comparing the phonetic geometry induced by different
supervision signals.


\subsection{Evaluation of Phonetic Representation Spaces}
\label{sec:lr-evaluation}

\subsubsection{Independent Phonetic Similarity Measures}
\label{sec:lr-distance}

Pronunciation accuracy evaluates whether a model selects the correct label but
does not establish how distances between different pronunciations are
represented internally. Independent phonetic distance measures provide an
alternative criterion. Li et al. introduced DIMSIM, which estimates Mandarin
phonetic similarity using learned representations of initials and finals
together with lexical-tone information \cite{ref_dimsim}. Its parameters were
informed by annotated phonetic relationships, providing a reference independent
of the training objectives considered in the present study. Such measures allow
hidden-state distances to be evaluated against graded pronunciation
relationships.

\subsubsection{Representational Similarity Analysis}
\label{sec:lr-rsa}

Representational similarity analysis compares the structure of representations
through pairwise dissimilarities \cite{ref_kriegeskorte2008}. In Chinese reading
research, this approach has been used to compare phonological dissimilarity with
patterns of neural activity \cite{ref_li_chinese2022}. Contextual language
representations introduce additional variation because the same linguistic item
may receive different hidden states across occurrences. Previous studies have
investigated context-sensitive similarity judgments and the geometry of
contextualized embeddings \cite{ref_huang2012,ref_ethayarajh2019}. These methods
provide a basis for examining relationships between pronunciations using
representations extracted from written character sequences.

\subsubsection{Comparison and Statistical Inference}
\label{sec:lr-inference}

The evaluation of representational geometry requires distinguishing systematic
relationships from variation introduced by the selected observations. Schütt et
al. examined statistical inference for representational similarity measures,
including resampling procedures that account for dependencies between
observations \cite{ref_schutt2023}. Phonetic discrimination studies have also
assessed learned representations through relative comparisons between speech
sounds \cite{ref_matusevych2023}. Together, these approaches motivate evaluating
both the correspondence between pairwise distances and the consistency of
relative similarity relationships.
